\section{Conclusão}

Este trabalho analisou três problemas clássicos de escalonamento: Machine
Scheduling, Job Shop Scheduling e Flow Shop Scheduling. As formulações foram
implementadas em Julia, com JuMP e o solver HiGHS, e avaliadas sobre 48
instâncias registradas nos arquivos de resultados: 31 de Machine Scheduling, 10
de Job Shop Scheduling e 7 de Flow Shop Scheduling. No conjunto consolidado,
foram obtidos 29 resultados com status ótimo, distribuídos entre os três
problemas, sem registros de erro de execução nos resultados computacionais
consolidados.

Os resultados indicaram comportamentos distintos entre os problemas. O Flow Shop
Scheduling apresentou a maior proporção de soluções ótimas, com 6 instâncias
ótimas em 7, além dos menores tempos médio e mediano. O Job Shop Scheduling
apresentou os maiores tempos médio e mediano, o que evidencia maior dificuldade
computacional no comportamento central observado. O Machine Scheduling, embora
tenha apenas uma máquina nas instâncias analisadas, concentrou a maior
quantidade absoluta de interrupções por limite de tempo e os maiores gaps,
especialmente nas instâncias com maior número de tarefas. Assim, considerando
conjuntamente tempo, gap e limites de tempo, os resultados apontam maior
dificuldade prática para Machine Scheduling no fechamento da otimalidade e para
Job Shop Scheduling no esforço computacional típico.

Os gráficos de Gantt complementaram a análise agregada ao mostrar a estrutura
das soluções selecionadas. Em Machine Scheduling, a programação da instância
\texttt{inst\_n05\_s01} concentrou toda a carga em uma única máquina, que
determinou diretamente o makespan. Em Job Shop Scheduling, a instância
\texttt{ft06} evidenciou o efeito das precedências entre operações e das esperas
causadas por conflitos de máquina, com algumas máquinas atuando como gargalos
mais relevantes. Em Flow Shop Scheduling, a instância
\texttt{problem\_3m\_10j} mostrou o fluxo das tarefas pela sequência comum de
máquinas, com formação de gargalo na máquina de maior carga e presença de
tempos ociosos decorrentes do encadeamento entre estágios.

A configuração adotada, com limite de tempo de 300 segundos e gap relativo de
1\%, foi adequada para obter soluções viáveis em todas as instâncias registradas
e para resolver parte expressiva do conjunto até a otimalidade segundo o
critério do solver. Contudo, essa configuração não foi suficiente para fechar o
gap em todas as instâncias, principalmente em Machine Scheduling e Job Shop
Scheduling. A interpretação dos resultados deve considerar essa limitação, bem
como o fato de que os conjuntos possuem tamanhos diferentes e que os resultados
dependem das formulações MILP diretas utilizadas.

Como trabalhos futuros, recomenda-se avaliar formulações alternativas,
estratégias de fortalecimento das restrições, diferentes valores de Big-M,
heurísticas de geração de soluções iniciais e outros limites de tempo ou
tolerâncias de gap. Também seria útil armazenar os horários completos de início
e término para todas as instâncias, permitindo validações estruturais e gráficos
de Gantt sem necessidade de nova resolução dos modelos. Essas extensões
permitiriam uma comparação mais ampla entre formulações e uma análise mais
detalhada do comportamento computacional dos três problemas.
